\documentclass[10pt,a4paper]{amsart}
\usepackage[margin=19mm]{geometry}
\usepackage{amsmath,amssymb,mathtools}
\usepackage[scr=rsfs,cal=euler]{mathalfa}
\usepackage{xfrac}
\usepackage[unicode,hidelinks,bookmarks=false]{hyperref}
\newcommand{\ie}{i.e.\ }
\newcommand{\cf}{cf.\ }
\newcommand{\resp}{respectively\ }
\newcommand{\Subsection}{\subsection}
\providecommand{\nobreakdash}{\nobreak\hbox{-}}
\setlength{\emergencystretch}{2em}
\multlinegap0pt
\usepackage{amssymb}
\usepackage{tensor}
\usepackage{mathtools}
\usepackage{dsfont}
\usepackage[shortlabels]{enumitem}
\newtheorem{lemma}{Lemma}
\title{A Multiplicative-Noise Mechanism for Variability Amplification under Radiative Forcing in an Arctic Energy-Balance Model}
\author{Gianmarco Del Sarto, Franco Flandoli, Marta Lenzi}
\date{Source: arXiv:2603.01617v1 (CC BY 4.0)}
\begin{document}
\maketitle

\noindent\emph{Source and licence.} A Multiplicative-Noise Mechanism for Variability Amplification under Radiative Forcing in an Arctic Energy-Balance Model. Version arXiv:2603.01617v1 (CC BY 4.0), distributed by arXiv under the Creative Commons Attribution 4.0 International licence, http://creativecommons.org/licenses/by/4.0/, as recorded in the arXiv licence metadata for this exact version and shown on the abstract page https://arxiv.org/abs/2603.01617v1. Full licence notice: \texttt{licenses/arxiv-2603.01617-CC-BY-4.0.txt}.\par\smallskip

\noindent\textit{Editorial scope.} Counted unit: the article's lemma lemma (source0.tex lines 1467--1477). The article's complete original proof of it is delivered with it (source0.tex lines 1478--1532, one \texttt{proof} environment of 55 lines). No mathematics is written, repaired or re-derived: the statement and the proof are the article's own lines, in the article's own notation, and no retained line is substituted or re-flowed.  No further source-local context is needed: every cross-reference of every retained line resolves inside the two delivered spans.  Nothing was omitted from any retained span, so the package carries no omission record.  No retained line contains a citation, so the wrapper carries no bibliography and no bibliography entry.  No retained line contains a figure environment, an image-inclusion command or drawing code, so no image is delivered and the package declares no asset dependency.  Editorial formatting only, in addition: an AMS \texttt{amsart} preamble that restates the article's own declarations and macros used by the retained lines, namely the environments \texttt{lemma} on separate counters, together with the standard packages those lines need.  The retained environments print in this document as Lemma 1 in order of appearance, whereas the article prints the same items with the numbering of its own full document.  The complete native source of the article as distributed in the arXiv e-print for this exact version is bundled unchanged as \texttt{sources/195567/source0.tex}.
\begin{lemma}
    Assume that $M^\lambda = A_\Delta -B^\lambda$ is negative definite and $y_0 = 0$. Then $ \Gamma_t^\lambda$ satisfies the matrix ODE
    \begin{equation*}
            \begin{split}
                 \frac{d}{dt} \Gamma_t^\lambda &= M^\lambda \Gamma_t^\lambda + \Gamma_t^\lambda (M^\lambda)^T
+ \mathbb{E}\!\left[\Sigma^\lambda(Y_t^\lambda)\,(\Sigma^\lambda(Y_t^\lambda))^T\right]\\
                 & = M ^\lambda \Gamma_t^\lambda + \Gamma_t^\lambda (M^\lambda)^T + \tau \, \sum_{k=1}^d \mathrm{diag}(l^k) (D^\lambda \Gamma_t^\lambda (D^\lambda)^T + f^\lambda (f^\lambda)^T) \mathrm{diag} (l^k),
            \end{split}
        \end{equation*}
        where $l^k$ is the $k$-th column of $L.$
\end{lemma}

\begin{proof}
    Assume for simplicity $\tau = 1$. We drop for the ease of notation the dependence on $\lambda$. Using It\^{o}-formula
    \begin{equation*}
        \begin{split}
            d (Y_t Y_t^T) &= (dY_t)(Y_t)^T + Y_t(dY_t)^T + (dY_t)(dY_t)^T  \\
            & = (MY_t dt + \Sigma  dW_t) (Y_t)^T + Y_t( MY_tdt+ \Sigma  dW_t)^T + (dY_t)(dY_t)^T.
        \end{split}
    \end{equation*}
    Note that, using It\^{o} isometry, we have
    \begin{equation*}
        \begin{split}
            (dY_t)(dY_t)^T =(MY_tdt+ \Sigma  dW_t) (MY_tdt+ \Sigma  dW_t)^T = \Sigma  dW_t (dW_t)^T \Sigma^T.
        \end{split}
    \end{equation*}
    Since $dW_t (dW_t)^T = I dt$, we obtain
    $$
    (dY_t)(dY_t)^T = \Sigma \Sigma^T dt.
    $$
    In conclusion we have
    $$
    d(Y_t Y_t^T) = [M Y_t(Y_t)^T  + Y_t (Y_t)^T M^T + \Sigma \Sigma^T]dt + \widetilde M_t,
    $$
    where $\widetilde M_t$ denotes a martingale. Taking the expected values, we deduce
    $$
    \frac{d}{dt} \Gamma_t = M  \Gamma_t + \Gamma_t  M^T + \mathbb{E}[\Sigma \Sigma^T].
    $$
    But in our setting $\Sigma = \mathrm{diag} (DY_t +f)L$, thus
    $$
    \Sigma \Sigma^T = \mathrm{diag} (DY_t +f) L L^T \mathrm{diag}( DY_t +f)^T.
    $$
    Since
    $$
    L L^T = \sum_{k =1 }^d l^k (l^k)^T,
    $$
    we obtain
    $$
    \Sigma \Sigma^T = \sum_{k = 1}^d \mathrm{diag} ( DY_t + f) l^k (l^k)^T \mathrm{diag} (DY_t +f)^T,
    $$
    Further, using the property
    $$
    \mathrm{diag} (a) b = \mathrm{diag}(b)a, \quad a,b \in \mathbb{R}^d,
    $$
    we have
    $$
    \Sigma \Sigma^T = \sum_{k =1 }^d \mathrm{diag} (l^k ) \left[(DY_t +f) (DY_t +f)^T \right] \mathrm{diag}(l^k).
    $$
    Since $\mathbb{E}[Y_t] = 0$, we have
    $$
    \mathbb{E}\left[(DY_t +f) (DY_t+f)^T \right] = D \Gamma_t  D^T + f f^T,
    $$
    and thus
    $$
    \mathbb{E}\left[\Sigma \Sigma^T \right] = \sum_{k = 1}^d \mathrm{diag}(l^k) (D \Gamma_t  D^T + f f^T) \mathrm{diag} (l^k).
    $$
\end{proof}

\end{document}
